\section{模型假设}
\begin{enumerate}[label=(\arabic*), leftmargin=2.5em, itemsep=0.3em]
\item 问题一为确定性问题：假设题目给出的光伏发电预测值即为当天的实际出力，\textbf{不考虑预测误差}；该假设\textbf{仅适用于问题一}。
\item 每个 $10\,$min 离散时段内，小区负载、光伏发电功率与外网电价\textbf{保持不变}——三者在该尺度上变化很小，故用时段平均值代表其功率水平。
\item 储能设备的充电效率与放电效率均为 \textbf{90\%}，即充电时只有 90\% 的电能存入储能，放电时需额外释放 $1/0.9$ 倍电量；除充放电损耗外，\textbf{不考虑自放电及其他能量损失}。
\item 研究时段内储能设备的容量（12000\,kWh）与最大充放电功率（5000\,kW）\textbf{保持稳定}，不随使用过程衰减。
\item \textbf{问题二至问题四中}，每日 0:00 制定计划时\textbf{仅使用该时刻之前可获得的历史数据}及题目给出的预报，不使用当日实际运行数据，以保证策略的\textbf{因果性}。
\item 微网\textbf{仅能从外部电网购电}，不考虑向外部电网反向售电；当光伏发电富余且储能无法继续接纳时，\textbf{允许弃光}。
\item \textbf{问题二}按“\textbf{照付不议}”方式结算计划购电量，计划一经确定\textbf{不再调整}，供电不足时按当时电价的 5 倍紧急购电；问题三、问题四的调整机制按题目规定的 \textbf{0.5 倍 / 1.5 倍}电价另行结算。
\item 日内实际运行中，储能设备可根据实际负载与光伏相对预测的偏差\textbf{实时调节充放电功率}，并始终满足\textbf{容量、功率与效率约束}。
\item \textbf{（问题三、问题四）}光伏预报与调度时段\textbf{口径一致}：发布时刻后第 $k$ 小时的预报表征该整点小时内 6 个 $10\,$min 区间的平均功率，实际值按同一口径统计。
\item \textbf{（问题四）}每日 0:00 制定计划时，当日逐时电价\textbf{已由日前市场公布}，即电价\textbf{可知但逐日、逐时段波动}；电价未知情形作为模型推广方向讨论。
\item 忽略预测计算、优化求解与控制指令传输造成的\textbf{时间延迟}，认为计划与调度指令可即时执行。
\end{enumerate}
